% Original CSC302 teaching adaptation of Heather Macbeth, The Mechanics of Proof.
% Source: https://hrmacbeth.github.io/math2001/05_Logic.html
% Source snapshot: e660f42b13ddcb6d12b52ba036d6bd071a0cfb9b (Lean 4.3).
% All teaching prose and tables are original; mathematical source locations below.
% Lean snippets are demonstrations in the textbook environment, not build-verified.
% CSC302 lecture slides adapted from Heather Macbeth, The Mechanics of Proof.
\documentclass[aspectratio=169,11pt]{beamer}
\usepackage[T1]{fontenc}
\usepackage{lmodern}
\usepackage{amsmath,amssymb,mathtools}
\usepackage{booktabs,array,listings,tikz}
\usetikzlibrary{arrows.meta,positioning,calc}
\definecolor{ProofNavy}{HTML}{18324F}
\definecolor{ProofTeal}{HTML}{007D82}
\definecolor{ProofPaper}{HTML}{F3F6F8}
\usetheme{default}
\setbeamercolor{normal text}{fg=ProofNavy,bg=white}
\setbeamercolor{structure}{fg=ProofTeal}
\setbeamercolor{frametitle}{fg=ProofNavy}
\setbeamercolor{block title}{fg=white,bg=ProofNavy}
\setbeamercolor{block body}{fg=ProofNavy,bg=ProofPaper}
\setbeamercolor{block title example}{fg=white,bg=ProofTeal}
\setbeamercolor{block body example}{fg=ProofNavy,bg=ProofPaper}
\setbeamerfont{frametitle}{size=\Large,series=\bfseries}
\setbeamertemplate{navigation symbols}{}
\setbeamertemplate{itemize item}{\small\textbullet}
\setbeamertemplate{footline}{%
  \leavevmode\hbox{%
  \begin{beamercolorbox}[wd=.86\paperwidth,ht=2.8ex,dp=1.2ex,leftskip=1.6em]{author in head/foot}%
    \scriptsize CSC302\quad\textbar\quad\insertshorttitle
  \end{beamercolorbox}%
  \begin{beamercolorbox}[wd=.14\paperwidth,ht=2.8ex,dp=1.2ex,rightskip=1.6em plus1fil]{date in head/foot}%
    \scriptsize\insertframenumber/\inserttotalframenumber
  \end{beamercolorbox}}}
\setbeamersize{text margin left=8mm,text margin right=8mm}
\lstdefinelanguage{Lean}{%
  morekeywords={example,theorem,lemma,def,by,calc,have,exact,apply,intro,intros,constructor,rcases,rfl,rw,ring,numbers,linarith,nlinarith,rel,extra,addarith,obtain,left,right,use,choose,dsimp,change,ext,fun,show,from,at,with,induction,case,decide,contradiction,push_neg,by_cases,by_contra,assumption,existsi,simp,simpa},
  sensitive=true,morecomment=[l]{--},morestring=[b]"}
\lstset{language=Lean,basicstyle=\ttfamily\small,keywordstyle=\color{ProofTeal}\bfseries,
  commentstyle=\color{gray},stringstyle=\color{ProofNavy},showstringspaces=false,
  columns=fullflexible,keepspaces=true,breaklines=true,frame=single,
  rulecolor=\color{ProofTeal!30},backgroundcolor=\color{ProofPaper},
  aboveskip=0.5em,belowskip=0.4em}
\newcommand{\N}{\mathbb{N}}
\newcommand{\Z}{\mathbb{Z}}
\newcommand{\Q}{\mathbb{Q}}
\newcommand{\R}{\mathbb{R}}
\newcommand{\coursechapter}[2]{%
  \title[Chapter #1: #2]{Chapter #1\\#2}%
  \subtitle{The Mechanics of Proof \textbar{} CSC302 lecture notes}%
  \author{Based on Heather Macbeth's textbook}%
  \date{}%
}
\newcommand{\sourcecredit}[1]{\par\vfill{\tiny Source: Heather Macbeth, \emph{The Mechanics of Proof}, #1.}}
\newcommand{\bookurl}{\url{https://hrmacbeth.github.io/math2001/}}
\hypersetup{colorlinks=true,urlcolor=ProofTeal,linkcolor=ProofNavy}

\coursechapter{5}{Logic}
\begin{document}
\begin{frame}\titlepage\end{frame}

\begin{frame}{What you should be able to do}
\begin{itemize}
\item Check a propositional formula with a truth table and prove a logical equivalence.
\item Track how quantifier order controls the dependence of witnesses.
\item Use excluded middle without first deciding which case is true.
\item Move negations through connectives and quantifiers correctly.
\item Turn a negative statement into a concrete proof task.
\end{itemize}
\begin{block}{Three levels of reasoning}
Truth tables analyze propositions. Quantifiers introduce objects. Mathematics supplies facts about those objects.
\end{block}
Lean fragments use the textbook's educational tactic environment.
\sourcecredit{Chapter 5, Sections 5.1--5.3}
\end{frame}

\section{5.1 Logical equivalence}
% Section 5.1, Examples 5.1.1--5.1.2 and Exercise 5.1.3.
\begin{frame}{5.1: The connectives and their truth conditions}
Let $P$ and $Q$ be propositions.
\begin{center}
\begin{tabular}{cc|ccccc}
\toprule
$P$&$Q$&$\neg P$&$P\land Q$&$P\lor Q$&$P\to Q$&$P\leftrightarrow Q$\\\midrule
T&T&F&T&T&T&T\\
T&F&F&F&T&F&F\\
F&T&T&F&T&T&F\\
F&F&T&F&F&T&T\\\bottomrule
\end{tabular}
\end{center}
\begin{block}{The one false implication case}
$P\to Q$ fails exactly when $P$ is true and $Q$ is false.
When $P$ is false, there is no instance of $P$ that violates the implication.
\end{block}
\sourcecredit{Section 5.1, Example 5.1.2 and Exercise 5.1.3}
\end{frame}

% Exercise 5.1.3. Table independently computed.
\begin{frame}{Build a compound truth table from the inside out}
Consider $P\leftrightarrow(\neg P\lor Q)$.
\begin{center}
\begin{tabular}{cc|c|c|c}
\toprule
$P$&$Q$&$\neg P$&$\neg P\lor Q$&$P\leftrightarrow(\neg P\lor Q)$\\\midrule
T&T&F&T&T\\
T&F&F&F&F\\
F&T&T&T&F\\
F&F&T&T&F\\\bottomrule
\end{tabular}
\end{center}
Compute $\neg P$ first, then the disjunction, then the equivalence.
The final column matches $P\land Q$.
\begin{block}{Check versus explanation}
For propositional formulas, matching truth tables establishes classical equivalence. A structural proof explains how to convert a proof in each direction.
\end{block}
\sourcecredit{Section 5.1, Exercise 5.1.3 and Example 5.1.5}
\end{frame}

% Example 5.1.4.
\begin{frame}[fragile]{Logical equivalence means mutual implication}
Two statements $A$ and $B$ are logically equivalent when $A\leftrightarrow B$.
For example, $(P\lor P)\leftrightarrow P$.
\begin{lstlisting}
-- Demonstration: every branch of the disjunction gives P.
example (P : Prop) : Iff (Or P P) P := by
  constructor
  . intro h
    obtain h1 | h2 := h
    . exact h1
    . exact h2
  . intro h
    left
    exact h
\end{lstlisting}
\sourcecredit{Section 5.1, Example 5.1.4; ASCII teaching adaptation}
\end{frame}

% Example 5.1.5.
\begin{frame}{Distributing a conjunction over a disjunction}
\[
P\land(Q\lor R)\quad\Longleftrightarrow\quad
(P\land Q)\lor(P\land R).
\]
\begin{columns}[T]
\begin{column}{.47\textwidth}
\begin{block}{Forward}
We have $P$ and either $Q$ or $R$.
\begin{itemize}
\item If $Q$, build $P\land Q$ and choose the left alternative.
\item If $R$, build $P\land R$ and choose the right alternative.
\end{itemize}
\end{block}
\end{column}
\begin{column}{.47\textwidth}
\begin{block}{Backward}
Split the given disjunction.
\begin{itemize}
\item From $P\land Q$, keep $P$ and introduce $Q\lor R$ on the left.
\item From $P\land R$, keep $P$ and introduce $Q\lor R$ on the right.
\end{itemize}
\end{block}
\end{column}
\end{columns}
\sourcecredit{Section 5.1, Example 5.1.5}
\end{frame}

% Example 5.1.6; converse counterexample is an original extension.
\begin{frame}{A common witness works for every input}
Let $X,Y$ be types and let $P(x,y)$ be a proposition for $x\in X$, $y\in Y$.
\[
(\exists x\in X,\ \forall y\in Y,\ P(x,y))
\quad\Longrightarrow\quad
(\forall y\in Y,\ \exists x\in X,\ P(x,y)).
\]
\begin{enumerate}
\item Obtain a witness $x_0$ that works for every $y$.
\item Fix an arbitrary $y$.
\item Use that same $x_0$; the universal hypothesis gives $P(x_0,y)$.
\end{enumerate}
The converse can fail. Over $\Z$, take $P(x,y)$ to mean $x>y$.
For each $y$, $x=y+1$ works; no one integer is greater than every integer.
\sourcecredit{Section 5.1, Example 5.1.6; converse counterexample is an extension}
\end{frame}

% Example 5.1.6.
\begin{frame}{Negating existence: two direct implications}
\[
\neg\exists x\in A,\ P(x)\quad\Longleftrightarrow\quad
\forall x\in A,\ \neg P(x).
\]
\begin{block}{Forward}
Assume no witness exists. For an arbitrary $a\in A$, if $P(a)$ held, $a$ would be a witness. Thus $\neg P(a)$.
\end{block}
\begin{block}{Backward}
Assume every $a\in A$ fails $P$. If a witness $w$ existed, we would have both $P(w)$ and $\neg P(w)$.
\end{block}
This equivalence also holds when $A$ is empty. No choice of a default element is needed.
\sourcecredit{Section 5.1, Example 5.1.6}
\end{frame}

% Exercises 5.1.7, idempotence, De Morgan, and transport of exists.
\begin{frame}{5.1 Practice: build both conversions}
\begin{enumerate}
\item Prove $(P\land P)\leftrightarrow P$.
\item Prove $\neg(P\lor Q)\leftrightarrow(\neg P\land\neg Q)$.
\item Suppose $\forall x\in A,\ P(x)\leftrightarrow Q(x)$. Prove
\[
(\exists x\in A,\ P(x))\leftrightarrow(\exists x\in A,\ Q(x)).
\]
\end{enumerate}
In item 3, must the witness change when you convert the statement?
\sourcecredit{Section 5.1, Exercises 5.1.7}
\end{frame}

\begin{frame}{5.1 Solution ideas}
\begin{enumerate}
\item Forward: keep one copy of $P$. Backward: use the proof of $P$ twice to build $P\land P$.
\item Forward: a proof of either $P$ or $Q$ would give the forbidden disjunction. Backward: split a hypothetical $P\lor Q$; each branch contradicts the corresponding negation.
\item Forward: obtain $w$ and $P(w)$, apply the forward implication at $w$, and keep $w$ as the witness for $Q$. Backward: use the reverse implication at the same witness.
\end{enumerate}
The common structure is to unpack the given statement and rebuild the target.
\sourcecredit{Section 5.1, Exercises 5.1.7}
\end{frame}

\section{5.2 The law of the excluded middle}
% Section 5.2, Examples 5.2.5 and 5.2.6.
\begin{frame}{5.2: A case split need not decide the proposition}
\begin{block}{Law of the excluded middle}
For every proposition $P$, we may reason by the two cases $P$ and $\neg P$.
\end{block}
To prove a target $Q$:
\begin{enumerate}
\item Assume $P$ and prove $Q$.
\item Assume $\neg P$ and prove $Q$.
\item Conclude $Q$, because the cases cover all possibilities.
\end{enumerate}
In Lean, \texttt{by\_cases h : P} creates these two branches.
This is a classical logical principle, not an algorithm for deciding arbitrary propositions.
\sourcecredit{Section 5.2, Examples 5.2.5 and 5.2.6}
\end{frame}

% Definition in Section 5.2 and Examples 5.2.1--5.2.4.
\begin{frame}{A predicate with easy endpoints and a difficult middle}
Following the textbook, call $k\in\N$ \emph{superpowered} when
\[
S(k):\quad\forall n\in\N,\quad k^{\,k^n}+1\text{ is prime}.
\]
Exponentiation is grouped as $k^{(k^n)}$; the convention is $k^0=1$.
\begin{center}
\begin{tabular}{cll}
\toprule
$k$ & Useful calculation & Conclusion\\\midrule
$0$ & $0^{0^0}+1=1$ & $\neg S(0)$\\
$1$ & $1^{1^n}+1=2$ for every $n$ & $S(1)$\\
$3$ & $3^{3^0}+1=4$ & $\neg S(3)$\\\bottomrule
\end{tabular}
\end{center}
For $k=2$, the first values are $3,5,17,257,65537$. A finite list of successes does not prove $S(2)$.
\sourcecredit{Section 5.2, Examples 5.2.1--5.2.4}
\end{frame}

% Example 5.2.5.
\begin{frame}{Prove a boundary exists without locating it first}
We know $S(1)$ and $\neg S(3)$. Prove
\[
\exists k\in\N,\quad S(k)\land\neg S(k+1).
\]
\begin{center}
\begin{tabular}{lll}
\toprule
Case & Choose & Why it works\\\midrule
$S(2)$ & $k=2$ & $S(2)$ and $\neg S(3)$\\[2mm]
$\neg S(2)$ & $k=1$ & $S(1)$ and $\neg S(2)$\\\bottomrule
\end{tabular}
\end{center}
Both cases produce a witness. This argument proves existence even before we determine which case holds.
\sourcecredit{Section 5.2, Example 5.2.5}
\end{frame}

\begin{frame}[fragile]{The same boundary argument in abstract Lean}
\begin{lstlisting}
-- Demonstration: the argument uses only S 1 and Not (S 3).
example (S : Nat -> Prop) (h1 : S 1) (h3 : Not (S 3)) :
    Exists (fun k => And (S k) (Not (S (k + 1)))) := by
  by_cases h2 : S 2
  . use 2
    constructor
    . exact h2
    . exact h3
  . use 1
    constructor
    . exact h1
    . exact h2
\end{lstlisting}
The logical structure does not depend on the arithmetic inside $S$.
\sourcecredit{Section 5.2, Example 5.2.5; abstract teaching adaptation}
\end{frame}

% Example 5.2.6 and Exercises 5.2.7 (contraposition).
\begin{frame}{Double negation and contraposition}
\begin{block}{Double negation elimination}
Given $\neg\neg P$, split into $P$ or $\neg P$. The first case proves $P$ directly. The second contradicts $\neg\neg P$. Thus $P$ follows classically.
\end{block}
\begin{block}{Contraposition}
\[
(P\to Q)\quad\Longleftrightarrow\quad(\neg Q\to\neg P).
\]
Forward: assuming $\neg Q$ and $P$, the implication supplies the forbidden $Q$.
Backward: assume $P$. If $Q$ failed, the contrapositive would give $\neg P$; hence $Q$ holds by classical reasoning.
\end{block}
The forward implication is direct; the reverse uses double negation elimination.
\sourcecredit{Section 5.2, Example 5.2.6 and Exercises 5.2.7}
\end{frame}

% Exercises 5.2.7, contraposition and Euler factorization.
\begin{frame}{5.2 Practice: identify the classical step}
\begin{enumerate}
\item From $\neg P\to\neg Q$ and $Q$, prove $P$. Show explicitly where excluded middle is used.
\item You are given the exact factorization
\[
2^{2^5}+1=4294967297=641\cdot6700417.
\]
Use it to select a concrete $k$ satisfying $S(k)\land\neg S(k+1)$, where $S$ is the superpowered predicate.
\end{enumerate}
Why are five prime values insufficient to establish a universal statement?
\sourcecredit{Section 5.2, Exercises 5.2.7}
\end{frame}

\begin{frame}{5.2 Solution ideas}
\begin{enumerate}
\item Split into $P$ and $\neg P$. In the first case, we are done. In the second, the given implication produces $\neg Q$, contradicting $Q$. Excluded middle licenses that case split.
\item The value at $n=5$ has a nontrivial factor $641$, so $\neg S(2)$. We already know $S(1)$; therefore choose $k=1$.
\end{enumerate}
\begin{block}{What a counterexample must provide}
To refute $\forall n,\ P(n)$, one failing input is enough. To prove it, checking a fixed finite sample is generally insufficient.
\end{block}
\sourcecredit{Section 5.2, Exercises 5.2.7}
\end{frame}

\section{5.3 Normal form for negations}
% Section 5.3, Example 5.3.1 and Table 5.1.
\begin{frame}{5.3: Push negations through propositional operations}
\begin{center}
\begin{tabular}{lll}
\toprule
Original & Equivalent form & Textbook Lean lemma\\\midrule
$\neg\neg P$ & $P$ & \texttt{not\_not}\\[2mm]
$\neg(P\lor Q)$ & $\neg P\land\neg Q$ & \texttt{not\_or}\\[2mm]
$\neg(P\land Q)$ & $\neg P\lor\neg Q$ & \texttt{not\_and\_or}\\[2mm]
$\neg(P\to Q)$ & $P\land\neg Q$ & \texttt{not\_imp}\\\bottomrule
\end{tabular}
\end{center}
\begin{block}{Interpret the last row}
An implication fails when its premise holds and its conclusion fails. Negating an implication does \emph{not} negate both sides.
\end{block}
These are classical equivalences; some directions require excluded middle.
\sourcecredit{Section 5.3, Table 5.1}
\end{frame}

% Table 5.1 and Examples 5.3.2--5.3.4.
\begin{frame}{Negating quantifiers changes who chooses the object}
\[
\neg\exists x\in A,\ P(x)\quad\Longleftrightarrow\quad
\forall x\in A,\ \neg P(x),
\]
\[
\neg\forall x\in A,\ P(x)\quad\Longleftrightarrow\quad
\exists x\in A,\ \neg P(x).
\]
\begin{itemize}
\item No witness works $\Rightarrow$ every candidate fails.
\item Not every candidate works $\Rightarrow$ there is a counterexample.
\end{itemize}
For ordered numbers, also use
\[
\neg(a<b)\Longleftrightarrow b\le a,\qquad
\neg(a\le b)\Longleftrightarrow b<a.
\]
The domain of each quantifier stays the same throughout the transformation.
\sourcecredit{Section 5.3, Table 5.1 and Example 5.3.3}
\end{frame}

% Example 5.3.2.
\begin{frame}{Normalize one connective at a time}
Start with
\[
\neg\bigl(\forall m\in\Z,\ m\ne2\Rightarrow\exists n\in\Z,\ n^2=m\bigr).
\]
\begin{enumerate}
\item Negate the universal quantifier:
\[
\exists m\in\Z,\ \neg\bigl(m\ne2\Rightarrow\exists n\in\Z,\ n^2=m\bigr).
\]
\item Negate the implication:
\[
\exists m\in\Z,\ m\ne2\land\neg\exists n\in\Z,\ n^2=m.
\]
\item Negate the existential quantifier:
\[
\exists m\in\Z,\ m\ne2\land\forall n\in\Z,\ n^2\ne m.
\]
\end{enumerate}
\sourcecredit{Section 5.3, Example 5.3.2}
\end{frame}

% Example 5.3.1.
\begin{frame}{Why negating ``and'' requires a case split}
Prove $\neg(P\land Q)\Rightarrow\neg P\lor\neg Q$.
\begin{enumerate}
\item Suppose $\neg(P\land Q)$.
\item By excluded middle, either $P$ or $\neg P$.
\item If $\neg P$, choose the left alternative.
\item If $P$, show $\neg Q$: an assumed $Q$ would combine with $P$ to give the forbidden $P\land Q$. Choose the right alternative.
\end{enumerate}
Conversely, if either $P$ or $Q$ is false, they cannot both hold.
\begin{block}{Lean connection}
Use \texttt{constructor}, then \texttt{by\_cases hP : P} in the forward direction; use \texttt{obtain} in the reverse direction.
\end{block}
\sourcecredit{Section 5.3, Example 5.3.1}
\end{frame}

% Example 5.3.3.
\begin{frame}{A chained inequality hides a conjunction}
The statement $n^2<m<(n+1)^2$ means
\[
n^2<m\ \land\ m<(n+1)^2.
\]
Therefore
\[
\neg\bigl(\forall n\in\Z,\ \exists m\in\Z,\ n^2<m<(n+1)^2\bigr)
\]
is equivalent to
\[
\exists n\in\Z,\ \forall m\in\Z,\quad
m\le n^2\ \lor\ (n+1)^2\le m.
\]
To \emph{prove} the resulting statement, choose $n=0$. Every integer $m$ satisfies $m\le0$ or $1\le m$.
\sourcecredit{Section 5.3, Example 5.3.3; witness proof is an extension}
\end{frame}

% Examples 5.3.4--5.3.5; Exercise 5.3.6 revisiting 4.5.1.
\begin{frame}[fragile]{Normalize first, then supply the mathematical idea}
\begin{lstlisting}
-- Demonstration in the textbook's tactic environment.
example : Not (forall x : Real, x ^ 2 >= x) := by
  push_neg
  use (1 / 2 : Real)
  numbers
\end{lstlisting}
After \texttt{push\_neg}, the goal asks for a real $x$ with $x^2<x$.
The tactic exposes the witness obligation; the choice $x=\tfrac12$ is the mathematical step.
\begin{block}{What normalization does}
It changes the form of the statement by equivalences. It does not automatically find the counterexample or prove the resulting inequality.
\end{block}
\sourcecredit{Section 5.3, Example 5.3.4 and Exercises 5.3.6}
\end{frame}

% Exercises 5.3.6, negation normalizations concerning products and upper bounds.
\begin{frame}{5.3 Practice: normalize, then prove or refute}
\begin{enumerate}
\item Move all negations inward:
\[
\neg\bigl(\forall a,b\in\Z,\ ab=1\Rightarrow(a=1\lor b=1)\bigr).
\]
Then decide whether the statement is true.
\item Normalize
\[
\neg\bigl(\exists x\in\R,\ \forall y\in\R,\ y\le x\bigr),
\]
and prove the normalized statement.
\end{enumerate}
Keep the variable domains visible while changing the logical symbols.
\sourcecredit{Section 5.3, Exercises 5.3.6}
\end{frame}

\begin{frame}{5.3 Solution ideas}
\begin{enumerate}
\item The normalized form is
\[
\exists a,b\in\Z,\quad ab=1\land a\ne1\land b\ne1.
\]
It is true: choose $a=b=-1$. Their product is $1$, and neither equals $1$.
\item The normalized form is
\[
\forall x\in\R,\ \exists y\in\R,\quad y>x.
\]
For an arbitrary $x$, choose $y=x+1$. Then $y>x$.
\end{enumerate}
A successful normalization turns the original wording into explicit inputs and witnesses.
\sourcecredit{Section 5.3, Exercises 5.3.6}
\end{frame}

% Example 5.3.5; natural-square separation proof.
\begin{frame}{One theorem, two useful proof presentations}
\[
\neg\exists n\in\N,\ n^2=2
\quad\Longleftrightarrow\quad
\forall n\in\N,\ n^2\ne2.
\]
\begin{block}{Proof of the normalized form}
Fix $n\in\N$. Either $n\le1$ or $n\ge2$.
\[
n\le1\Rightarrow n^2\le1<2,
\qquad
n\ge2\Rightarrow n^2\ge4>2.
\]
Either strict inequality implies $n^2\ne2$.
\end{block}
The numerical argument is unchanged from Chapter 4, but normalization lets us begin with an arbitrary input rather than a hypothetical witness.
\sourcecredit{Section 5.3, Example 5.3.5}
\end{frame}

\begin{frame}{Review: four questions before manipulating a statement}
\begin{enumerate}
\item What is the outermost connective or quantifier?
\item Am I proving an implication, an equivalence, or the truth of a particular formula?
\item Which variables may a witness depend on?
\item Does this direction use excluded middle?
\end{enumerate}
\begin{block}{A reliable workflow}
Write the domains $\longrightarrow$ expose hidden conjunctions $\longrightarrow$ normalize one layer at a time $\longrightarrow$ identify the resulting proof obligations.
\end{block}
A truth table checks propositional structure; a quantified claim also needs reasoning about its objects.
\sourcecredit{Chapter 5, Sections 5.1--5.3}
\end{frame}

\begin{frame}{Sources and further work}
Heather Macbeth, \emph{The Mechanics of Proof}, Chapter 5.
\begin{itemize}
\item 5.1: logical equivalence, truth tables, and quantifiers.
\item 5.2: excluded middle, the superpowered predicate, and contraposition.
\item 5.3: Table 5.1 and worked negation normalizations.
\end{itemize}
\medskip
\href{https://hrmacbeth.github.io/math2001/05_Logic.html}{Read Chapter 5 and practice the remaining logical equivalences.}
\medskip

Teaching prose, tables, and added counterexamples are original. Source comments record example locations. Lean snippets use the textbook's educational tactics.
\end{frame}
\end{document}
